\documentclass[11pt,a4paper]{article}

\usepackage[margin=1in]{geometry}
\usepackage{amsmath,amssymb,amsthm}
\usepackage{enumitem}
\usepackage{hyperref}
\usepackage{url}

\hypersetup{
  colorlinks=true,
  linkcolor=blue,
  citecolor=blue,
  urlcolor=blue
}

\newtheorem{theorem}{Theorem}[section]
\newtheorem{lemma}[theorem]{Lemma}
\newtheorem{definition}[theorem]{Definition}

\newcommand{\Tensor}{\operatorname{Tensor}}
\newcommand{\bandpart}{\operatorname{band\_part}}
\newcommand{\rowop}{\operatorname{row}}
\newcommand{\concat}{\mathbin{+\mkern-4mu+}}
\newcommand{\XEq}{\approx}
\newcommand{\SEq}{\simeq}
\newcommand{\Rstar}{\mathbb{R}^*}

\title{Formal Verification of KV-Cache Correctness\\for Sliding-Window Causal Attention}
\author{LiZhi Zhou\\[0.3em]\small Project artifact: \url{https://github.com/math-proof/lemma}}
\date{}

\begin{document}
\maketitle

\begin{abstract}
Production GPT-style decoders rely on a key--value (KV) cache to avoid recomputing past keys and values at every autoregressive step.
Implementations in Hugging Face Transformers, vLLM, and related systems concatenate cached tensors with the new token and attend over a sliding causal window, but this incremental update is validated by tests rather than proved as a tensor identity.
We state and machine-check such an identity in the \texttt{lemma} library: if every prefix output $Z(n)$ equals sliding-window masked scaled dot-product attention, then $Z(n{+}1)$ equals $Z(n)$ concatenated with one new attention row computed from the cached window plus the new key and value.
The proof is checked in Lean~4 (2026-08-19) and was first stated interactively in SymPy (2024-02-28).
Masking uses hyperreal $-\infty$ via a \texttt{band\_part} mask (library-defined from \texttt{tril}/\texttt{triu}/\texttt{masked\_fill}, following TensorFlow's \texttt{tf.linalg.band\_part}~\cite{tensorflow_band_part}), and equality is the library's total hyperreal closeness relation $\XEq$, designed to mimic \texttt{torch.isclose} on $\Rstar$.
We give the formal statement, a proof sketch, the lemma dependency structure, and limitations of what is and is not claimed.
Interactive visualizations and source files are public on \url{http://www.lemma.cn/} and GitHub~\cite{lemma_repo}.
\end{abstract}

\section{Introduction}
\label{sec:intro}

GPT-2, GPT-3, GPT-4, and every Hugging Face \texttt{AutoModelForCausalLM} decode one token at a time~\cite{radford2019language}.
At each step a self-attention layer maps hidden states to queries, keys, and values $Q,K,V$.
For the new token at position $n$, causal attention is
\[
\operatorname{softmax}\!\left(\frac{q_n K_{:n+1}^{\top}}{\sqrt{d_z}}+M\right)V_{:n+1},
\]
where $K_{:n+1}$ and $V_{:n+1}$ collect keys and values of all tokens seen so far and $M$ hides future positions~\cite{vaswani2017attention}.
The linear maps producing $K$ and $V$ do not depend on later tokens, so the first $n$ key/value rows never change; recomputing them at every generate step is quadratic in prefix length and wasted work.

A \emph{KV cache} stores past keys and values and, at the next step, concatenates only the new pair:
\[
K_{:n+1}=K_{:n}\concat[k_n],\qquad V_{:n+1}=V_{:n}\concat[v_n].
\]
Hugging Face Transformers implements this as \texttt{use\_cache=True} (the default in \texttt{generate})~\cite{huggingface_generate}.
Tensors travel as \texttt{past\_key\_values}; the default container is \texttt{DynamicCache}~\cite{huggingface_dynamic_cache}, one $(K,V)$ pair per layer with shape \texttt{[batch, heads, seq\_len, head\_dim]}.
Internally the attention module concatenates current $k,v$ with the cache and attends with a mask of length \texttt{past\_kv\_length + 1}~\cite{huggingface_kv_cache}.
Prefill runs the full prompt once; decode then feeds one new \texttt{input\_id} per step plus the cache.
Sliding-window models keep only the last $\ell$ pairs, matching our Lean statement; vanilla GPT keeps the whole prefix, matching the SymPy special case.

That is an implementation.
The contribution of this paper is the missing algebraic identity: if $Z(n)$ is already masked attention of the prefix, then ``concatenate cache with new $k,v$ and take the last row'' is exactly $Z(n{+}1)$, not a different algorithm.
We prove this as a machine-checked lemma in a \texttt{Tensor} type semantically equivalent to \texttt{torch.Tensor}~\cite{lemma_repo}, with causal masking $(\bandpart-1)\cdot\infty$ over hyperreals $\Rstar$ so masked positions are identically zero after $\exp$, rather than a floating-point stand-in such as $-10^9$.

Public interactive statements appear on lemma.cn in SymPy~\cite{lemma_kv_cache_sympy} and Lean~4~\cite{lemma_kv_cache_lean}; source files are
\texttt{All\_Eq\_DotSoftmaxAdd\_DivDot\_T.lean} and \texttt{kv\_cache.py} on GitHub~\cite{lemma_repo}.

\section{Related work}
\label{sec:related}

\textbf{Runtime systems.}
PyTorch, Hugging Face Transformers~\cite{huggingface_kv_cache}, vLLM~\cite{vllm}, and SGLang~\cite{sglang} implement KV caches and sliding windows as production systems.
Correctness is tested, not proved as a tensor equation.

\textbf{Formal libraries.}
Mathlib~\cite{mathlib_hyperreal} has hyperreal analysis but no \texttt{torch.Tensor} calculus, no \texttt{band\_part}, and no softmax-attention lemmas matching PyTorch semantics.
OrthoCache machine-checks Parseval for the Walsh--Hadamard transform and a total-variation bound for spectral eviction~\cite{orthocache}; it does not prove that an incremental cached row equals full masked attention.
TorchLean / BatchInvariantInference~\cite{torchlean} specifies schedule-explicit CUDA attention and batch invariance---a refinement of kernels against a denotation, not this library's \texttt{Tensor} identity with \texttt{band\_part} and hyperreal $\infty$.
The \texttt{fak} project reports bit-identical cache tests ($\max|\Delta|=0$) in Go; those are empirical witnesses, not Lean/SymPy theorems.

\textbf{Hyperreal closeness.}
Classical nonstandard analysis~\cite{keisler1976elementary} and Mathlib define ``infinitely close'' by $a-b$ infinitesimal, equivalently $\operatorname{st}(a)=\operatorname{st}(b)$ when both are finite.
Infinite hyperreals have no standard part, so that relation is not total on $\Rstar$.
\texttt{torch.isclose}~\cite{pytorch_isclose} is IEEE floating-point code: $|x-y|\le \mathrm{rtol}\cdot|y|+\mathrm{atol}$; non-finite values are close iff equal.
It is not a theorem and not hyperreal.
Our library's $\XEq$ relation~\cite{lemma_xeq} fills this gap: a machine-checked, total relation on $\Rstar$ mimicking \texttt{torch.isclose} with infinitesimal relative tolerance.
The KV-cache lemma is among the first large tensor identities using it: $\exp(-\infty)\XEq 0$, and masked softmax is $\XEq$ the windowed stack rather than definitionally equal in $\mathbb{R}$.

\section{Preliminaries}
\label{sec:prelim}

\subsection{Tensors and notation}

\texttt{Tensor $\mathbb{R}$ s} is this library's formalization of a real tensor of shape \texttt{s}, with operations matching \texttt{torch.Tensor}: matrix multiply (\texttt{@}), transpose $(\cdot)^{\top}$, \texttt{softmax}, \texttt{tril}, \texttt{triu}, \texttt{masked\_fill}, slicing \texttt{[start:stop]}, stacking $[i<n]\,f(i)$, and concatenation $\concat$.
Lifting to \texttt{Tensor $\Rstar$ s} interprets entries in the hyperreals.
Shape-cast equality is written $\SEq$ (\texttt{SEq}).

\paragraph{\texttt{band\_part} (not a native PyTorch API).}
PyTorch has no \texttt{band\_part}.
This library names the TensorFlow operation \texttt{tf.linalg.band\_part}~\cite{tensorflow_band_part} and defines it from PyTorch primitives~\cite{lemma_repo}:
\begin{verbatim}
# matches Tensor.band_part in special.lean (d=1: no extra masked_fill)
def band_part(X, l, u, d=1):
    Y = (X.tril(u).triu(-l))
    if d == 1:
        return Y
    delta = j - i   # diagonal index on the last two axes
    return Y.masked_fill((delta + l) % d != 0, 0)
\end{verbatim}
\noindent
Here \texttt{j - i} is the diagonal index $j-i$ on the last two axes (as in \texttt{sympy/matrices/expressions/special.lean}).
Entry $(i,j)$ is kept iff $j-i\in[-l,u]$ and, when $d>1$, iff $d\mid(j-i+l)$ (dilated band; default $d=1$).
Lemma \texttt{Tensor.BandPart.eq.Stack\_BoolIn\_Icc} proves this equals the boolean indicator of $j-i\in[-l,u]$ for $d=1$.
We write $\bandpart(l,u)$ for \texttt{band\_part(ones(n,n), l, u)}.
\emph{Do not confuse} the band parameters $(l,u)$ with the \emph{sliding-window width} $\ell$ below: the causal window mask uses $\bandpart(\ell-1,0)$, which keeps exactly the last $\ell$ keys per row ($j-i\in[1-\ell,0]$).

Fix sliding-window width $\ell\ge 1$ (the number of past keys each query may attend to), embedding size $d_z$, and sequences
\[
Q,K,V:\mathbb{N}\to\Tensor\mathbb{R}[d_z],
\]
together with prefix outputs $Z(n)\in\Tensor\mathbb{R}[n,d_z]$.
Write $Q_{:n}=[i<n]\,Q(i)$ and likewise for $K,V$.
\emph{Notation:} $\ell$ is the mathematical symbol for this window width; Lean and code use the same name \texttt{l} ($\ell=\texttt{l}$).
The \texttt{band\_part} mask in the theorems is \texttt{band\_part(..., l-1, 0)}, \emph{not} \texttt{band\_part(..., l, 0)}---the first argument to \texttt{band\_part} is a band limit, not the window width.

\subsection{Hyperreal closeness \texorpdfstring{$\XEq$}{approx}}

\begin{definition}[\texorpdfstring{$\XEq$}{approx} on \texorpdfstring{$\Rstar$}{R*}~\cite{lemma_xeq}]
\label{def:xeq}
For $a,b\in\Rstar$,
\[
a\XEq b
\quad\Longleftrightarrow\quad
\frac{|a-b|}{|a|+|b|+1}\to 0,
\]
i.e.\ the relative gap is infinitesimal.
On tensors, $\XEq$ is the pointwise lift of this scalar relation.
\end{definition}

This is \emph{not} Mathlib's infinitely-close relation ($a-b$ infinitesimal).
It mimics \texttt{torch.isclose}~\cite{pytorch_isclose}:
absolute tolerance is effectively $0$ (infinitesimal absolute error implies $\XEq$, but not conversely);
relative tolerance is the displayed quotient, with $+1$ in the denominator for stability near zero.
Two infinities may differ by an infinite amount and still be $\XEq$, which a naive $|a-b|$ test forbids.

Throughout the proof, $\exp(A+(\Xi-1)\infty)\XEq e^{A}\Xi$ is false as an identity of reals and false as NSA infinitely-close when both sides are infinite; it is true as $\XEq$.

\subsection{Hyperreal band-part mask}

\begin{lemma}[Mask exponential~\cite{lemma_gpt}]
\label{lem:mask}
With $\Xi_{ij}=\mathbf{1}_{[-\ell+1,0]}(j-i)$,
\[
\exp\bigl(A_{ij}+(\Xi_{ij}-1)\infty\bigr)\XEq e^{A_{ij}}\Xi_{ij}.
\]
On the support of $\Xi$ the infinite term vanishes; off the support the logit is $-\infty$ and the exponential is infinitesimally $0$.
Softmax of that matrix is the sliding-window causal softmax without a finite clip constant (Figure~\ref{fig:mask}).
\end{lemma}

\begin{figure}[ht]
\centering
\[
\Xi_{ij}=\mathbf{1}_{[-\ell+1,0]}(j-i),\qquad
\exp\bigl(A_{ij}+(\Xi_{ij}-1)\infty\bigr)\XEq e^{A_{ij}}\Xi_{ij}.
\]
\caption{Hyperreal \texttt{band\_part} mask: forbidden logits go to $-\infty$; $\exp$ yields infinitesimal zero off the sliding window.}
\label{fig:mask}
\end{figure}

\subsection{Stack decomposition (lemma \texttt{gpt})}

Lemma \texttt{gpt}~\cite{lemma_gpt} unfolds batched sliding-window masked attention into one row per token---the row-wise form used inside a KV cache.
Throughout, $\ell\ge 1$ is the \emph{sliding-window width}: query at row $i$ attends to keys with indices in $[(i+1-\ell)_+,i+1)$ (at most $\ell$ keys, causal).
The Lean proof uses parameter \texttt{l} $=\ell$ and mask \texttt{band\_part(l-1, 0)}.

\begin{lemma}[Stack decomposition (sliding window $\ell$)~\cite{lemma_gpt}]
\label{lem:gpt}
Fix sliding-window width $\ell\ge 1$.
For $Q,K,V\in\Tensor\mathbb{R}[n,d_z]$,
\begin{align*}
&\operatorname{softmax}\!\left(
\frac{QK^{\top}}{\sqrt{d_z}}
+\bigl((1_{[n,n]}).\bandpart(\ell-1,0)-1\bigr)\infty
\right)V \\
&\qquad\XEq\;
[i<n]\;
\operatorname{softmax}\!\left(\frac{Q[i]\,K[(i+1-\ell):(i+1)]^{\top}}{\sqrt{d_z}}\right)
V[(i+1-\ell):(i+1)].
\end{align*}
The left-hand side is batched masked attention with a causal band of width~$\ell$ (mask $\bandpart(\ell-1,0)$).
The right-hand side is the same computation row by row: each row $i$ attends only to its length-$\ell$ key/value window---exactly the slice a KV cache concatenates and reuses at decode step~$i$.
When $\ell\ge n$ the window is the full prefix (vanilla GPT causal attention).
\end{lemma}

\section{Main result}
\label{sec:main}

\begin{theorem}[KV-cache increment~\cite{lemma_kv_cache_lean}]
\label{thm:kv-cache}
Assume $\ell\ge 1$ and for every $n$,
\begin{equation}
\label{eq:hyp}
Z(n)\;\XEq\;
\operatorname{softmax}\!\left(
\frac{Q_{:n}K_{:n}^{\top}}{\sqrt{d_z}}
+\bigl((1_{[n,n]}).\bandpart(\ell-1,0)-1\bigr)\cdot\infty
\right)
V_{:n}.
\end{equation}
The mask is $\bandpart(\ell-1,0)$, i.e.\ in PyTorch (with \texttt{l} $=\ell$ the sliding-window width):
\begin{verbatim}
# l == ell in the paper (sliding-window width; Lean name l)
torch.ones(n, n).tril(0).triu(-(l - 1))
\end{verbatim}
\noindent
equivalently \texttt{band\_part(torch.ones(n, n), l - 1, 0)} in this library.
Entry $(i,j)$ is kept iff $j-i\in[1-\ell,0]$, i.e.\ causal attention restricted to the last $\ell$ keys.

Let
\begin{align*}
K_w&=K_{:n}[(n+1-\ell):n],&
V_w&=V_{:n}[(n+1-\ell):n],\\
\rowop&=
\operatorname{softmax}\!\left(\frac{Q(n)\,(K_w^{\top}\concat K(n)^{\top})}{\sqrt{d_z}}\right)
(V_w\concat[V(n)]).
\end{align*}
Then
\begin{equation}
\label{eq:conc}
Z(n+1)\;\XEq\; Z(n)\concat[\rowop].
\end{equation}
\end{theorem}

Thus one new query attends only to the cached window plus the new key/value; the previous $n$ output rows are reused unchanged.

The SymPy theorem~\cite{lemma_kv_cache_sympy} is the full-causal special case \texttt{BandPart[n, 0]}, with
\[
\operatorname{softmax}\!\left(\frac{Q[n]\,[K[:n]^{\top}\mid K[n]]}{\sqrt{d_z}}\right)[V[:n]\mid V[n]],
\]
i.e.\ the window is the entire prefix.
When $\ell>n+1$ in Lean, slice $[(n+1-\ell):n]$ is the whole prefix and the SymPy theorem is recovered.

\section{Proof sketch}
\label{sec:proof}

The proof is fully checked in Lean~4 in a single file importing the lemmas listed in Section~\ref{sec:deps}.
Figure~\ref{fig:reduction} summarizes the reduction.

\begin{figure}[ht]
\centering
\[
\begin{array}{c}
Z(n)\;\XEq\;
\operatorname{softmax}\!\left(
\dfrac{Q_{:n}K_{:n}^{\top}}{\sqrt{d_{z}}}
+\bigl((1).\bandpart(\ell-1,0)-1\bigr)\infty
\right)V_{:n}
\quad(\forall n)
\\[1.1em]
\Big\downarrow\ \scriptstyle\text{instantiate at }n+1,\ \text{Lemma~\ref{lem:gpt}}
\\[0.9em]
Z(n+1)\;\XEq\;
[i<n+1]\;
\operatorname{softmax}\!\left(
\dfrac{Q[i]\,K[(i+1-\ell):(i+1)]^{\top}}{\sqrt{d_{z}}}
\right)
V[(i+1-\ell):(i+1)]
\\[1.1em]
\Big\downarrow\ \scriptstyle [i<n+1]f(i)=[i<n]f(i)\concat[i<1]f(n)
\\[0.9em]
\begin{array}{c@{\qquad}c}
[i<n]\,f(i) & [f(n)] \\[0.5em]
\Big\downarrow\ \scriptstyle h(n)+\text{Lemma~\ref{lem:gpt}} & \Big\downarrow\ \scriptstyle\text{slice}=\text{window}\concat\text{new token} \\[0.5em]
Z(n) &
\operatorname{softmax}\!\left(\dfrac{Q(n)\,(K_{w}^{\top}\concat K(n)^{\top})}{\sqrt{d_{z}}}\right)(V_{w}\concat[V(n)])
\end{array}
\\[1.6em]
Z(n+1)\;\XEq\; Z(n)\concat[\rowop]
\end{array}
\]
\caption{Proof reduction for Theorem~\ref{thm:kv-cache}.}
\label{fig:reduction}
\end{figure}

\paragraph{Step 1: instantiate and unfold.}
Apply hypothesis~\eqref{eq:hyp} at $n+1$ and Lemma~\ref{lem:gpt} to express $Z(n+1)$ as a stack of windowed row attentions.

\paragraph{Step 2: split the stack.}
By \texttt{Tensor.Stack.eq.AppendStackS},
$[i<n+1]f(i)=[i<n]f(i)\concat[i<1]f(n)$ for the row function $f$ from Lemma~\ref{lem:gpt}.
Congruence of $\XEq$ under $\concat$ is \texttt{Tensor.XEqAppendS.of.XEq.XEq}.

\paragraph{Step 3: identify the prefix block.}
The block $[i<n]f(i)$ equals $Z(n)$ by the hypothesis at $n$ and Lemma~\ref{lem:gpt}, after proving window slices of $K_{:n+1}$ and $K_{:n}$ agree below $n$ via \texttt{Tensor.SEqGetSliceSStack.of.LengthSlice} and \texttt{Tensor.GetSliceStack.as.Stack\_UFnAdd.of.Eq}.

\paragraph{Step 4: identify the last row.}
The singleton $[f(n)]$ is attention of $Q(n)$ against
$K_{:n+1}[(n+1-\ell):(n+1)]\SEq K_w\concat[K(n)]$ (and likewise for $V$).
Transposing concatenated row-blocks is concatenation of transposes up to shape cast (\texttt{Tensor.TAppend.as.AppendTS}).
Composing with softmax-scale-matmul yields $\rowop$, giving~\eqref{eq:conc}.

\section{Formalization artifact}
\label{sec:artifact}

Interactive visualizations of the finished statement:
\begin{itemize}[noitemsep]
  \item SymPy: \url{http://www.lemma.cn/py/?module=Tensor.GetSlice.eq.Append_DotSoftmaxDivDot_Append.of.All_Eq_DotSoftmaxAdd_DivDot_T.kv_cache}
  \item Lean~4: \url{http://www.lemma.cn/lean/?module=Tensor.GetSlice.eq.Append_DotSoftmaxDivDot_Append.of.All_Eq_DotSoftmaxAdd_DivDot_T}
\end{itemize}
The SymPy page displays \texttt{apply}/\texttt{prove} for prefix equality implying concatenation with \texttt{row}.
The Lean page displays \texttt{@[main] private lemma kv\_cache} with \texttt{[NeZero (l : $\mathbb{N}$)]}, universal hypothesis on $Z(n)$, the \texttt{let} chain \texttt{Kn, Vn, Kw, Vw, KT, kT, row}, and conclusion $Z(n+1)\XEq Z(n)\concat[\rowop]$.

Source: \url{https://github.com/math-proof/lemma} (Lean~4 on \texttt{main}, SymPy on \texttt{master}).

\subsection{Lemma dependencies}
\label{sec:deps}

Figure~\ref{fig:deps} lists major imported lemmas (private helpers in the same file omitted).

\begin{figure}[ht]
\small
\begin{verbatim}
Tensor.GetSlice.eq...All_Eq_DotSoftmaxAdd_DivDot_T  (kv_cache)
|-- Tensor.DotSoftmaxAdd_Mul_Infty.eq.Stack_DotSoftmaxDivDot_T  (gpt)
|   |-- Tensor.DotSoftmaxAdd_Mul_Infty.eq.Stack_DotSoftmax
|   |   |-- Tensor.BandPart.eq.Stack_BoolIn_Icc
|   |   |-- Tensor.ExpAdd_MulInfty.eq.Mul_Stack_Bool
|   |   `-- Tensor.Softmax.eq.DivExp_KeepdimSumExp
|   |-- Tensor.GetDot_TGetSlice.as.Dot_Get
|   `-- Tensor.GetSliceGetDiv.eq.DivGetSliceGet
|-- Tensor.Stack.eq.AppendStackS
|-- Tensor.XEqAppendS.of.XEq.XEq
|-- Tensor.GetSliceStack.as.Stack_UFnAdd.of.Eq
|-- Tensor.SEqGetSliceSStack.of.LengthSlice
|-- Tensor.SEqAppendS.of.SEq.SEq
`-- Tensor.TAppend.as.AppendTS
\end{verbatim}
\caption{Import graph for the Lean proof of Theorem~\ref{thm:kv-cache}.}
\label{fig:deps}
\end{figure}

\section{SymPy versus Lean~4}
\label{sec:sympy-lean}

The SymPy lemma (2024-02-28)~\cite{lemma_kv_cache_sympy} is the full-causal increment with \texttt{Equal} rather than $\XEq$.
Its proof applies the Python \texttt{gpt} unfolding, substitutes $n\mapsto n+1$, splits the stack, and rewrites last-row slices as \texttt{SEq\_Append}.

The Lean~4 lemma (2026-08-19)~\cite{lemma_kv_cache_lean} keeps parameter $\ell\neq 0$, slices $[(n+1-\ell):n]$, works in $\Rstar$ with $\XEq$, and carries \texttt{cast}s required by dependent transpose shapes.
A batched causal variant exists on the SymPy side at \texttt{Tensor.Eq.of.Eq.gpt.kv\_cache.batched}.

\section{Discussion and limitations}
\label{sec:discussion}

The statement is an implication, not an if-and-only-if:
\emph{given} $\forall n.\; Z(n)\XEq$ masked-attention$(Q_{:n},K_{:n},V_{:n};\ell)$,
\emph{imply} $Z(n+1)\XEq Z(n)\concat[\rowop]$.

It does not construct $Q,K,V$ from a neural network, does not address multi-head or grouped-query layouts, and does not prove CUDA-level bit identity.
It does prove that, in this \texttt{Tensor} calculus, the KV-cache update is the same mathematical object as one more step of sliding-window masked attention.

The lemma does not bound approximation error from dropping tokens outside the window (that is already encoded in the mask).
It does not replace OrthoCache's spectral eviction theorems~\cite{orthocache} or TorchLean's kernel-refinement theorems~\cite{torchlean}.
It is the complementary statement: if prefixes are exact masked attention, the cached one-row update is exact masked attention of length $n+1$.

\section{Conclusion}
\label{sec:conclusion}

We have presented a machine-checked identity linking GPT-style KV-cache updates to incremental sliding-window masked attention, together with a hyperreal closeness relation suitable for masked softmax with $-\infty$.
The result is public on lemma.cn and GitHub~\cite{lemma_repo}, with proofs checked in Lean~4 and an earlier SymPy formulation.
Future work includes multi-head layouts, batched Lean formalization, and connecting $\XEq$ to floating-point refinement theorems.

\bibliographystyle{plain}
\bibliography{refs}

\end{document}
